% ==========================================================================
% GENERATED by scripts/split_paper.py. Do not edit this file.
%
% Source: imported_papers/probe-norm/content/body.tex, content/supplement.tex
% Part 3: Materials and Methods
%
% Edits belong upstream, in the published paper. To change what the
% chapter includes, change imports/bounded.toml and re-run
% the splitter.
% ==========================================================================
\section{Materials and Methods}
\label{sec:methods}
\subsection{Norm drift and state probes}
\label{sec:background}
\label{sec:formal_setup}
Let the world model's encoder $\phi$ map an observation $\mathbf o_t$ as $\z_t=\phi(\mathbf o_t)\in\mathbb R^D$. Given actions $\mathbf a_{t+k}$, the dynamics module $F_\theta$ produces a rollout
\begin{equation}
\hat{\z}_t=\z_t,\qquad
\hat{\z}_{t+k+1}=F_\theta(\hat{\z}_{t+k},\mathbf{a}_{t+k}).
\label{eq:recurrence}
\end{equation}
Here $\z_{t+k}$ is the encoding of the observed future frame and $\hat\z_{t+k}$ is its prediction. Training supervises steps $k\le H$. Evaluation may extend to $T>H$.

A state probe $r:\mathbb R^D\to\mathbb R^{d_y}$ is fitted on reference encodings with physical states $\mathbf y$ as targets and then applied to predictions. Its mean squared error over $n$ rollouts and $d_y$ state coordinates is
\begin{equation}
E_r(T)=\frac{1}{n d_y}\sum_{i=1}^{n}
\bigl\|r(\hat{\z}^{(i)}_{t+T})-\mathbf{y}^{(i)}_{t+T}\bigr\|^2 .
\label{eq:eval_error}
\end{equation}
We call $E_r$ the \emph{raw error} and measure how it changes when prediction norms are matched to those of reference encodings.

\subsubsection{Where the norm enters the error}
\label{sec:split}
We can rewrite a nonzero prediction as $\hat\z=\hat c\,\hat{\mathbf u}$, where $\hat c=\|\hat\z\|$ is its norm and $\hat{\mathbf u}=\hat\z/\|\hat\z\|$ is its unit direction. Let $\rho(\hat{\mathbf u})$ denote its \emph{reference radius}, the norm that encodings in direction $\hat{\mathbf u}$ have. When the encoder ends in a normalization, we compute the radius from that normalization, which gives the layer-normalization radius for our trained models and the exact SimNorm radius for TD-MPC2 (\supsec~\ref{app:radius}). Otherwise we take the norm of the reference encoding nearest in angle.

Re-projection replaces $\hat\z$ with $\tilde\z=\rho(\hat{\mathbf u})\hat{\mathbf u}$. The error of a single prediction, $r(\hat\z)-\mathbf y$, is the term inside the sum of Equation~(\ref{eq:eval_error}). Adding and subtracting the probe output at the reference radius splits it into two parts:
\begin{equation}
\begin{aligned}
r(\hat{\z})-\mathbf y &= r(\hat{\z})-\mathbf y-r(\tilde\z)+r(\tilde\z) \\
&=\underbrace{\bigl[r(\hat{\z})-r(\tilde\z)\bigr]}_{\text{change due to the norm}}+\underbrace{\bigl[r(\tilde\z)-\mathbf y\bigr]}_{\text{residual after rescaling}} .
\end{aligned}
\label{eq:probe_composition}
\end{equation}
The first term is the change caused by the norm. The second is the error remaining at the reference radius. It contains the error of the predicted direction together with any error the probe makes on directions unlike the encodings it was fitted on, and re-projection cannot reduce it.

\subsubsection{When the norm effect is large}
\label{sec:loss}
\label{sec:arch_constraint}
\label{sec:readout}
For a linear probe $r(\z)=W\z+\mathbf b$~\citep{alain2016understanding,Hewitt2019ControlTasks}, the norm effect factors into
\begin{equation}
r(\hat\z)-r(\tilde\z)
=\underbrace{\bigl(\hat c-\rho(\hat{\mathbf u})\bigr)}_{\text{norm mismatch}}\quad
\underbrace{W\hat{\mathbf u}}_{\text{scale sensitivity}} .
\label{eq:linear_norm_effect}
\end{equation}
The score therefore changes only when the predicted norm differs from its reference radius and the probe is sensitive to scale, that is, when $W\hat{\mathbf u}$ is not zero. Three factors decide how large these two terms are.
\begin{itemize}
\item \textbf{Training loss.} Cosine loss ignores the norm of the vector it compares. Squared error regulates the norm, but only at supervised steps. Neither guarantees matching norms throughout a rollout longer than the training horizon ($T>H$).
\item \textbf{Output constraints.} Explicit normalization can prevent norm mismatch at every step. SimNorm, for example, splits the latent into groups and applies a softmax within each group~\citep{Hansen2023}, so used on both encoder and dynamics it sets the norm of every prediction from its direction. A looser bound only limits growth and need not place predictions at the reference radius.
\item \textbf{Probe response.} Linear and ReLU probes give increasingly large outputs as the input norm grows~\citep{Hein2019ReLU,Xu2021Extrapolate}. A probe with bounded hidden activations, such as tanh, has bounded output. Normalizing the probe input removes scale sensitivity, but may discard information carried by norm.
\end{itemize}

Unless the design sets the norm of every prediction from its direction, as SimNorm on both encoder and dynamics does, or the probe normalizes its input, the size of the effect has to be measured.

\subsubsection{Measuring the effect of the norm}
\label{sec:reprojection}
At each horizon $T$, we evaluate the same fitted probe on raw and re-projected predictions, without changing the recurrent rollout. The re-projected error $E_r^{\mathrm{rep}}$ and its ratio to the raw error $E_r$ are
\begin{equation}
E_r^{\mathrm{rep}}(T)
=\frac{1}{n d_y}\sum_{i=1}^{n}
\bigl\|r(\tilde\z^{(i)}_{t+T})-\mathbf y^{(i)}_{t+T}\bigr\|^2,
\qquad
q_r(T)=\frac{E_r^{\mathrm{rep}}(T)}{E_r(T)} .
\label{eq:reprojection_factor}
\end{equation}
The \emph{re-projection factor} $q_r(T)$ is defined when $E_r(T)>0$. Values below one indicate that rescaling reduces error. For example, $q_r=0.1$ corresponds to a $90\%$ reduction. Values above one indicate increased error. The factor measures how much of the raw error the norm causes, while $E_r^{\mathrm{rep}}$ measures the error that remains under the same probe.

Re-projection can also discard information. The norm can itself carry state information, and under squared-error training, averaging uncertain future directions can shorten a prediction, so restoring its norm may increase error. Four checks test whether a re-projected error reflects state information.

\subsubsection{Checks}
\label{sec:checks}
\begin{itemize}
\item \textbf{Support check.} Re-projection restores the norm, but the probe is reliable on the re-projected prediction only if its direction lies among the encodings it was fitted on. Far from every encoding, the probe extrapolates, and its error no longer measures the model. We measure the angle of each prediction to the nearest reference encoding and call the prediction \emph{supported} when this angle lies below the 99th percentile of the same angles for held-out target encodings.
\item \textbf{Radius check.} We fit two probes on reference encodings, one on the direction alone and one on direction and norm. The gain in held-out $R^2$ shows how much state the norm carries that the direction lacks, and that re-projection would remove.
\item \textbf{Probe agreement.} A ranking of two models can depend on the choice of probe. We therefore order each pair with five probes (Section~\ref{sec:instruments}): the re-projected and the normalized probe, which ignore the norm, and the probe fitted on predictions, the neighbor probe and target retrieval, which do not extrapolate beyond the data they were fitted on.
\item \textbf{Null controls.} A useful prediction must beat predictors that need no learned dynamics.
\begin{itemize}
\item \emph{Persistence} applies the probe to the initial encoding, as if the state did not change, and the \emph{reference mean} always predicts the mean state of the reference data.
\item Random directions with the prediction norms, re-projected in the same way, show how much rescaling lowers the error when the direction carries no state.
\item Three null rollouts, with actions shuffled across episodes, with starts swapped across episodes, and with the untrained module, test whether the score reflects the episode's actions, its initial state, or anything learned.
\item A \emph{start-encoding probe}, fitted to predict the future state directly from the initial encoding, tests whether the rollout adds information beyond the initial encoding.
\end{itemize}
\end{itemize}

\subsubsection{The protocol}
\label{sec:decision}
The protocol applies these steps in order at every evaluated horizon, with the rollout and the fitted probe fixed (Table~\ref{tab:terms}). The probe is first checked on encodings, including the radius check. We then use the re-projection factor to choose between the raw and the re-projected error, the null controls to test whether that error reflects learned dynamics, and the support check and probe agreement to test whether a ranking holds. A ranking is \emph{confirmed} when it passes every step, \emph{rejected} when its models fail the null controls, and \emph{inconclusive} otherwise. \supsec~\ref{app:code} gives code for each step.

\begin{table}[tb]
\caption{The protocol at one horizon, with the rollout and the probe fixed. Step 3 compares medians over runs.}
\label{tab:terms}
\centering\footnotesize
\begin{tabular}{@{}l>{\raggedright\arraybackslash}p{0.25\linewidth}>{\raggedright\arraybackslash}p{0.27\linewidth}>{\raggedright\arraybackslash}p{0.3\linewidth}@{}}
\toprule
Step & Question & Quantity & Decision \\
\midrule
1. Probe on encodings & Is the probe accurate on encodings, and does the norm carry state? & Held-out $R^2$, scale sensitivity, radius gain & $R^2$ above $0.5$ required; radius gain above $0.05$: use the raw error \\
2. Re-projection & How much of the error does the norm cause? & Factor $q_r$, beside the probe's error on the encoded targets & $q_r$ below $0.9$: use the re-projected error; otherwise the raw error \\
3. Null controls & Does the prediction beat what needs no learned dynamics? & Re-projected error against persistence, reference mean, start-encoding probe, random directions and null rollouts & Any of them as good: model rejected, and so is a ranking that puts it first or pairs two rejected models \\
4. Support and probes & Does the probe extrapolate, and does a ranking depend on the probe? & Share of predictions supported, orderings under five probes & Re-projected and normalized probes agree, both models at least half supported and no probe reverses: confirmed; otherwise inconclusive \\
\bottomrule
\end{tabular}
\end{table}

\subsection{Experimental setup}
\label{sec:variants}
\label{sec:suite}
We test re-projection and the protocol in three settings of increasing realism: synthetic latents whose state content we control, world models we train ourselves while varying how their architecture, loss and constraints act on the norm, and eleven published world models.

\subsubsection{Synthetic latents}
Synthetic latents let us set the state that a latent encodes and its norm independently, so we know what a correct probe should report. Section~\ref{sec:res_suite} describes the constructions, and \supsec~\ref{app:synthetic} gives their details.

\subsubsection{Trained world models}
\label{sec:protocol}
The norm of a prediction depends on three design choices: the architecture of the dynamics module, the training loss, and constraints on the latent. We train world models on pendulum, cartpole and acrobot swing-up~\citep{Tassa2018} and vary each choice. Every model is trained to predict $H=6$ steps and evaluated on held-out episodes with random actions, five and twenty steps ahead. Reference episodes fit the probes, and separate scoring episodes measure prediction error (\supsec~\ref{app:impl}).

To test the architecture, we train four dynamics modules with the same cosine loss and data, ten runs per module and task, which gives $120$ base models. Non-Linear, a residual network, and Linear, a linear layer with additive action inputs, leave the norm free, so we expect their norms to drift. Linear-Penalized adds a penalty on the linear map's spectral radius, which discourages growth without preventing it. Orthogonal constrains the linear map to be orthogonal, so it preserves the norm before the action is added, and we expect little drift. Each model learns its own encoder with representation regularization~\citep{Bardes2022}, and the encoder ends in layer normalization, which gives the reference radius analytically (\supsec~\ref{app:radius}).

To test the loss, we replace the cosine loss with squared error, which penalizes norm mismatch at supervised steps, so we expect drift only beyond the training horizon. To test constraints, we rescale every predicted latent to the initial norm, apply SimNorm to encoder and dynamics, or remove the encoder's final normalization. The first two should remove the drift, and the last should keep it. Each of these four variants has five runs per module and task.

To test the probe where actions matter and all modules share one latent space, the shared-encoder variant trains only the dynamics of all four modules on the frozen encoder of one Non-Linear run per task, under random torques held for four frames, with ten runs per module and task, evaluated up to fifty steps ahead.

To test rankings whose true answer is known, we build models of known accuracy on the encoders of the thirty Non-Linear base models. Such a model predicts, for each scoring episode, the encoding of a real frame from another episode whose state lies at a chosen neighbor rank from the true state five steps ahead, so a lower rank gives a more accurate model. We let the more accurate model drift furthest from the reference radius, which should mislead the ordinary probe, and give every model ten seeds (\supsec~\ref{app:known}).

\subsubsection{Published world models}
\label{sec:lewm}
The trained models test one design choice at a time. To test whether norm drift distorts evaluations in practice, we apply the protocol to eleven published world models on sixteen tasks: TCRL~\citep{zhao2023simplified}, TD-MPC and TD-MPC2~\citep{hansen2022temporal,Hansen2023}, LeWorldModel~\citep{maes2026leworldmodel}, Semigroup-JEPA~\citep{Liu2026SemigroupJEPA}, plain and consistent Koopman autoencoders~\citep{Azencot2020}, KSL~\citep{mcinroe2021learning}, PlayVirtual~\citep{yu2021playvirtual}, and the MuZero, EfficientZero and Sampled EfficientZero checkpoints of LightZero~\citep{Schrittwieser2020,ye2021mastering,hubert2021learning,niu2023lightzero}. We use released checkpoints or the authors' code and settings (\supsec~\ref{app:published}), and train six of the models ourselves (\suptab~\ref{tab:fidelity}). For Koopman autoencoders and Semigroup-JEPA, we keep the authors' own decoder or probe, so that we test their evaluation as published. For the other models, we fit the same state probe as for our trained models, so that every model is evaluated with the same probe.

\subsubsection{Probes}
\label{sec:instruments}
Our main probe, the \emph{ordinary probe}, is a multilayer perceptron with one hidden layer of $64$ ReLU units, fitted on reference encodings with three random seeds (\supsec~\ref{app:encoder}). The \emph{re-projected probe} is the same fitted probe applied to re-projected predictions. To test whether a ranking depends on how the latents are read, we compare it with four other probes that either ignore the norm or cannot extrapolate: the \emph{normalized probe}, fitted on unit directions; a probe fitted on predictions; the \emph{neighbor probe}, which returns the mean state of the nearest reference encodings; and \emph{target retrieval}, the rank of the correct target among all encoded targets (\supsec~\ref{app:alternative_readouts}).

\subsubsection{Statistical analysis}
\label{sec:stats}
The base models and the shared-encoder variant have ten independent runs for each of the four modules (Non-Linear, Linear, Linear-Penalized and Orthogonal) and three tasks (pendulum, cartpole and acrobot), and the squared-error, fixed-radius, SimNorm and no-final-normalization variants have five. Errors compound multiplicatively over a rollout, so small differences between runs grow to orders of magnitude at long horizons, and we report medians over runs. To compare two models, we pair every run of one with every run of the other and take the share of pairs in which the first has the lower error. The first model is better when the $95$\% bootstrap interval of this share lies above one half.

%%BEGIN_RESULTS
